\section{Inverse Jobs and Complete Filtering Chains}
\label{sec:inverse-chain}
Forward analysis is only one use of a resumable transform. The comparison
campaign also evaluates independently released inverse transforms and complete
frequency-domain filtering chains. These are benchmark applications of the
reusable FFT classes; neither visualization module produces audio or invokes
these paths. Their schedules retain bulk buffering operations, so the bounded
operation contract of Section~\ref{sec:one-hop} applies only to the complete
analysis implementation. This distinction lets the measurements expose the
remaining concentration of work in the inverse and filtering controls.

\subsection{Normalized inverse jobs}
An inverse job accepts all $N$ complex bins in natural order and produces all
$N$ complex time samples,
\begin{equation}
 z[m]=\frac{1}{N}\sum_{k=0}^{N-1}Z[k]e^{j2\pi km/N}
     =\frac{1}{N}\overline{\operatorname{FFT}(\overline Z)[m]}.
 \label{eq:inverse-contract}
\end{equation}
The spectrum need not be Hermitian. No window is applied or inverted, and the
factor $1/N$ belongs to every adapter's measured work. In
\code{OnTheFlyIFFT<T>}, buffering copies and permutes the input and conjugates
the retained coefficients. Each subsequent step executes one forward
butterfly or conjugates and normalizes one output. Thus the transform has
$B_c(N)+N$ resumable steps, with $B_c(N)=(N/2)\log_2N$.
Buffering remains a timed $O(N)$ pass outside that step count.

The streaming job fixture releases a new spectrum on sample $t_r=rH$,
independently of any forward analyzer. After the transform, the harness copies
all $N$ results to its canonical destination, giving
\begin{equation}
 W_{\mathrm{inv}}=B_c(N)+2N
 \label{eq:inverse-work}
\end{equation}
scheduled units in the first-party control. The immediate variant executes
them on release; the balanced variant uses Equation~\eqref{eq:balanced} with
$W=W_{\mathrm{inv}}$. Both pay the same bulk input preparation at release.
Their release-to-publication delays are respectively $d=0$ and $d=H-1$
samples. A native-library job executes its transform, required conversion,
normalization, and canonical stores immediately. These jobs publish $N$
complex outputs every $H$ calls; they do not define an audio stream with
$N=H$, perform overlap-add, or inherit the analyzer's spectrum-center age.

\subsection{Overlap-save processing and delivery}
The complete-chain fixture follows conventional FFT convolution
principles~\cite{wefers2015}. Let $v[n]$ be the common complex input stream,
zero before startup, and let $g[\ell]$ be a causal FIR of length $L$.
At endpoint $t_r$, the selected frame is
$u_r[m]=v[t_r-N+1+m]$. With $G[k]$ the $N$-point transform of the zero-padded
filter, the chain computes
\begin{equation}
 c_r=\operatorname{IFFT}\!\left(\operatorname{FFT}(u_r)\,G\right),
 \qquad b_r[j]=c_r[N-H+j],\quad 0\le j<H.
 \label{eq:chain-contract}
\end{equation}
The product is binwise, and both transforms are full complex transforms.
The last $H$ samples agree with linear convolution when $N-H\ge L-1$:
each retained index then has every required input within the selected frame,
so none of its filter terms wraps circularly. Discarding the first $N-H$
outputs is overlap-save extraction. No analysis or synthesis window is used.

The two fixtures are identity and
$g=[1/2,-1/4,1/8]$. Both require $H\le N-2$ and execute the same path,
including all $N$ spectral multiplications; the identity does not bypass the
transforms. The three-tap filter is a correctness and application fixture,
not a claim that FFT convolution is economical for such a short filter.
This experiment evaluates neither a windowed STFT reconstruction nor
partitioned long-filter convolution.

The first-party chain schedules forward butterflies, $N$ multiplications,
one inverse-buffer operation, inverse butterflies and normalization, and $H$
output stores. Its accounting is therefore
\begin{equation}
 W_{\mathrm{chain}}=2B_c(N)+2N+H+1.
 \label{eq:chain-work}
\end{equation}
The inverse-buffer operation charged as one unit contains $O(N)$ work.
At frame release, an additional $O(N)$ ring-to-frame copy and forward
buffer/permutation pass execute outside the quota. Consequently,
$\ceil{W_{\mathrm{chain}}/H}$ bounds scheduled units per call, but does not
bound their primitive operations by the analysis scheduler's formula.
The native-library chains execute the same functional pipeline immediately,
including layout conversion and normalized inverse output. All chains include
continuous delivery and a common dependent checksum that keeps every delivered
sample observable to the optimizer.

Publication exposes a block at $p_r=t_r+d$. Playback starts in that call and
delivers $b_r[j]$ at $p_r+j$. Since this sample corresponds to convolution
output index $t_r-H+1+j$, the input-to-delivery delay is
\begin{equation}
 D_{\mathrm{chain}}=H-1+d
 =\begin{cases}
 H-1,&\text{immediate},\\
 2H-2,&\text{balanced},
 \end{cases}
 \label{eq:chain-delay}
\end{equation}
in samples, or $D_{\mathrm{chain}}/f_s$ seconds. This is the implementation's
buffering delay relative to the causal convolution output; any filter phase
delay is separate. Successive blocks join without gaps after startup because
each contains $H$ outputs and publications are $H$ samples apart.
Device buffering and callback-end visibility add delay beyond this sample
index model.

\subsection{Independent output checks}
Alternating sparse non-Hermitian spectra have known inverse sequences of the
form $a+b e^{j2\pi 7m/N}+c e^{-j2\pi 3m/N}$, allowing every inverse output
to be checked without reusing a tested transform. For the chains, direct
time-domain identity/FIR evaluation checks every delivered sample and every
published block, including startup zeros and retained-input wraparound.
These synthesis checks require
$|a-r|\le\epsilon\max(1,|r|)$ for each complex output, with
$\epsilon=2\times10^{-5}$ for float and $10^{-10}$ for double; non-finite
values fail. This pointwise absolute-or-relative budget differs from the
processed-spectrum norm policy used for analysis. Correct fixture output and
exact logical delivery do not by themselves demonstrate audio-device
deadline compliance.
